% Complete cut factorization and both universal reductions for Kamp III.2, Lemma 7.

\section{Kamp's reductions across an interior cut}
\label{sec:annex-p09-kamp-iii2-lemma7}

\begin{annexitem}{annex:p9.kamp.cut.factorization}{thm:p09-kamp-cut-factorization}{Every point cut, gap cut and exact residual length}
The source splits its predicate occurrence sequence into a prefix and
suffix, then completes the two words at their common cut. The following
proof supplies the realization argument behind those operations and
checks every permitted length
\cite[ch.~III, pp.~64--66, 81--82]{Kamp1968}.

\begin{proof}
Fix an arbitrary nonempty strict linear structure \(\mathcal M\)
and a word \(W=(A_0,B_1,A_1,\ldots,B_n,A_n)\), with \(n\geq1\).
Keep the predicate interpretations fixed throughout. Write \(a<d<b\)
for the domain points assigned to the three endpoints \(r,s,t\).
The basic-pattern realization equivalence from
\cref{annex:p9.kamp.iii2.lemma1} identifies the original pattern
at \((a,b)\) with a strictly increasing realization
\[
                         a=a_0<a_1<\cdots<a_n=b.
\]
It also applies to each of the two cut words at its ordered endpoints.

\paragraph{The finite cuts are exhaustive.}
If \(d\) equals one of the \(a_i\), then \(a<d<b\)
forces \(1\leq i\leq n-1\); thus it is an interior point
cut. There is at most one such index because the sequence is strict.
If no equality holds, consider the nonempty finite set
\(\{i:a_i<d\}\). It contains 0 and excludes \(n\), so it
has a greatest member \(j\leq n-1\). By maximality and absence
of equality, \(a_j<d<a_{j+1}\). This is the gap cut with
\(k=j+1\in\{1,\ldots,n\}\). Thus every realization meets
exactly one of the listed positional cases. This argument takes a
maximum of a finite set of indices; it assumes no completeness of
the domain order.

\paragraph{Splitting at an interior point.}
Suppose \(d=a_k\), where \(1\leq k\leq n-1\). The prefix
\[
                         (a_0,\ldots,a_k)
\]
realizes \(W_{(\mathrm p,k)}^L\) from \(a\) to \(d\).
For its point positions \(i=0,\ldots,k\), the original sequence
supplies \(A_i(a_i)\). For every gap index \(i=1,\ldots,k\),
its open interval is the unchanged \((a_{i-1},a_i)\) and has
the unchanged predicate \(B_i\) throughout.
The suffix \((a_k,\ldots,a_n)\) realizes
\(W_{(\mathrm p,k)}^R\) from \(d\) to \(b\): its point
labels are \(A_k,\ldots,A_n\), and its gap labels are
\(B_{k+1},\ldots,B_n\), at exactly their original positions.
Both words require the common point label \(A_k(d)\), which the
original realization supplies. Each side has at least one interval
because \(k\) is strictly between 0 and \(n\).

\paragraph{Splitting inside a gap.}
Suppose \(a_{k-1}<d<a_k\), where \(1\leq k\leq n\).
Use
\[
             (a_0,\ldots,a_{k-1},d),\qquad
             (d,a_k,\ldots,a_n)
\]
as the two realizing sequences. Every old point on either side retains
its original \(A_i\) predicate. Every old gap other than the
one containing \(d\) remains unchanged with its original predicate.
The old condition
\((a_{k-1},a_k)\subseteq B_k^{\mathcal M}\) supplies all
three requirements at the cut:
\[
 B_k(d),\qquad
 (a_{k-1},d)\subseteq B_k^{\mathcal M},\qquad
 (d,a_k)\subseteq B_k^{\mathcal M}.
\]
These are exactly the two copies of \(B_k\) at the end of
\(W_{(\mathrm g,k)}^L\) and the two copies at the beginning
of \(W_{(\mathrm g,k)}^R\). The left word has \(k\)
intervals, and the right has \(n-k+1\), both positive. Hence
each split gives the indicated pair \(L_c(a,d),R_c(d,b)\).
Together with exhaustion, this proves the forward factorization.

\paragraph{Joining two realizations at a point cut.}
Conversely, suppose \(L_c(a,d)\land R_c(d,b)\) holds for
\(c=(\mathrm p,k)\). Choose realizations of the two cut words:
\[
       a=l_0<\cdots<l_k=d,
       \qquad d=r_k<r_{k+1}<\cdots<r_n=b.
\]
The indices on the right are chosen to match their original point
positions. Identify the common point \(l_k=r_k=d\), and take
\[
                         (l_0,\ldots,l_k,r_{k+1},\ldots,r_n).
\]
This sequence is strict: each constituent was strict, and its only
join has \(l_k=d<r_{k+1}\). It has \(n+1\) points and
endpoints \(a,b\). The left point labels supply
\(A_0,\ldots,A_k\); the right supply
\(A_{k+1},\ldots,A_n\), with the common \(A_k(d)\)
also supplied on the right. All gaps up to index \(k\) are
left gaps with \(B_1,\ldots,B_k\), and all later gaps are
right gaps with \(B_{k+1},\ldots,B_n\). No new gap is
created. The sequence therefore realizes the original word.

\paragraph{Joining two realizations through a gap cut.}
Suppose instead that \(c=(\mathrm g,k)\). Choose realizations
\[
       a=l_0<\cdots<l_{k-1}<d,
       \qquad d<r_k<r_{k+1}<\cdots<r_n=b.
\]
The left sequence has \(d\) as its final point and the right
sequence has \(d\) as its first point. Their endpoint labels
both assert \(B_k(d)\). Their adjoining gap conditions give
\[
       (l_{k-1},d)\subseteq B_k^{\mathcal M},\qquad
       (d,r_k)\subseteq B_k^{\mathcal M}.
\]
Delete the common cut point and join the remaining lists:
\[
                         (l_0,\ldots,l_{k-1},r_k,\ldots,r_n).
\]
It is strictly increasing because \(l_{k-1}<d<r_k\), and
it has exactly \(n+1\) points with endpoints \(a,b\).
All retained points have their original labels \(A_i\).
Gaps of indices less than \(k\) retain the left realization's
conditions, and those of indices greater than \(k\) retain
the right realization's conditions.

For the only combined gap, take \(z\) with
\(l_{k-1}<z<r_k\). Linearity gives the exhaustive cases
\(z<d\), \(z=d\), and \(d<z\). The left gap supplies
\(B_k(z)\) in the first case, the cut endpoint label in the
second, and the right gap in the third. Hence
\((l_{k-1},r_k)\subseteq B_k^{\mathcal M}\), as required.
In particular the test at \(d\) cannot be omitted: it is an
interior point of the recovered gap, even if the two old adjoining
gaps were empty. The joined sequence realizes \(W\), completing
the reverse factorization.

\paragraph{The exact bounds on both sides.}
A point cut \((\mathrm p,k)\) gives lengths
\[
              \ell(W_c^L)=2k+1,\qquad
              \ell(W_c^R)=2(n-k)+1.
\]
Because \(1\leq k\leq n-1\), both are at most \(2n-1\)
and hence strictly below \(\ell(W)=2n+1\).
A gap cut \((\mathrm g,k)\) gives
\[
              \ell(W_c^L)=2k+1,\qquad
              \ell(W_c^R)=2(n-k+1)+1.
\]
Both are at most \(2n+1\). Equality for the left occurs exactly
at \(k=n\), and equality for the right occurs exactly at
\(k=1\). Thus each complete side list has exactly one candidate
of the original length, and none longer. Restricting a list to any
subset of its indices preserves the stated at-most-one bound.
The claim concerns occurrences in the constructed lists: no semantic
deduplication is needed to make the bound true.

\paragraph{Boundary and syntactic cases.}
At \(n=1\) the point-cut list is empty and the sole gap cut
gives \((A_0,B_1,B_1)\) and \((B_1,B_1,A_1)\). The preceding
split and join arguments still apply, and both lengths are three.
At the first gap, the left sequence has just its original initial
point and the cut; at the last gap, the right sequence has just
the cut and the original final point. There are no omitted exterior
cuts because the given point satisfies \(a<d<b\).

Repeated predicate letters have one fixed interpretation wherever they
occur. Every use above was tied to a specific word position, so a
letter appearing in both point and gap positions does not alter the
choice of cut type. Empty open gaps impose their usual vacuous universal
condition. A gap cut nevertheless requires the actual point \(d\)
inside that old gap; a point cut requires it among the internal
witnesses. No step creates a new point. All formula constructors have
fresh bound variables, and their free variables are their two endpoints.
The realization equivalence therefore applies at precisely the guarded
assignments used here. This proves the asserted \(\mathsf{Lin}\)
identity and all its length statements.
\end{proof}
\end{annexitem}

\begin{annexitem}{annex:p9.kamp.iii2.lemma7.left}{thm:p09-kamp-iii2-lemma7-left}{The left-interior truth-pattern construction and its converse}
The finite truth-pattern step is Kamp's construction of
\(\sigma'_\alpha,\sigma_\alpha,\rho_\alpha\) and its final
display~(23). We give the missing cut justification through
\cref{thm:p09-kamp-cut-factorization}, and retain all Boolean
alternatives and the original length restriction
\cite[ch.~III, pp.~80--83]{Kamp1968}.

\begin{proof}
Fix a nonempty strict linear structure \(\mathcal M\), an arbitrary
interpretation of the predicates, and an assignment to every free
variable of the formulas in the statement. Let the assigned values of
\(x,y\) be \(a,b\). The truth of \(\psi_1\) below is
evaluated after assigning its displayed outer witness to \(u\).
Within the universal matrix, \(\psi_2\) is additionally evaluated
at the current value of \(v\). Every other free parameter keeps
its original value.

\paragraph{Finite complete truth patterns on the fixed side.}
For each function \(\alpha:\mathcal C_n\to\{0,1\}\), let
\[
 \Omega_\alpha(u,y)\defeq
 \bigwedge_{c\in\mathcal C_n}
 \begin{cases}
 R_c(u,y),&\alpha(c)=1,\\
 \neg R_c(u,y),&\alpha(c)=0.
 \end{cases}
\]
Here the cases select the formula placed at that position in the
conjunction. For any assignment to \(u,y\), exactly one such
truth pattern holds: assign 1 to precisely the true \(R_c\).
Every conjunct of that pattern holds by definition. A different
pattern disagrees at some index and has a false conjunct there.
Consequently \(\bigvee_\alpha\Omega_\alpha\) is valid.
This is the source's \(\sigma'_\alpha\), before its negative
conjuncts are removed. These formulas have no free occurrence of
\(v\); their truth is fixed during the universal quantification.

\paragraph{From the original formula to a single aggregate alternative.}
Assume the original formula is true. Its existential quantifier
supplies \(d\) with \(a<d<b\), with \(\psi_1\) true,
and with
\[
 \text{for every }z\in T\text{ such that }a<z<d,
       \quad \psi_2\text{ is true at }v=z
                      \text{ or }\phi(z,b)\text{ holds}.
 \tag{1}
\]
Define the one subset
\[
                     J=\{c\in\mathcal C_n:R_c(d,b)\text{ holds}\}.
\]
The complete fixed-side truth pattern selects this \(J\), and
all its positive members give \(\Sigma_J^L(d,b)\).
If \(J\) is empty, that conjunction is the stated tautology,
which also holds.

Take any \(z\) with \(a<z<d\). If \(\psi_2\) is
true at this value of \(v\), the required new disjunction is
true. Otherwise (1) gives \(\phi(z,b)\). The inequalities
\(z<d<b\) permit the cut factorization, which supplies some
\(c\in\mathcal C_n\) with
\(L_c(z,d)\land R_c(d,b)\). The second conjunct puts this
index in \(J\), and the first proves \(\rho_J^L(z,d)\).
Thus the new universal matrix holds for every such \(z\).
Outside that open interval the implication's guard is false.
The same witness \(d\) therefore satisfies the alternative
indexed by \(J\), and hence the complete output disjunction.

This step chooses the subset from the fixed-side truths once, before
considering individual \(z\). The cut index itself may vary
with \(z\). The residual disjunction allows this variation without
moving a choice of cut outside the universal quantifier incorrectly.
If \(J\) is empty, no point with false \(\psi_2\) can occur
in the tested interval: factorization at such a point would supply
a member of \(J\). Therefore the empty alternative is valid
with its false residual exactly when the unchanged \(\psi_2\)
supplies all required tests.

\paragraph{From any positive aggregate alternative to the original.}
Suppose an output alternative, indexed by some subset \(J\), is
true. Choose its existential witness \(d\). Then \(a<d<b\),
\(\psi_1\) holds, and \(R_c(d,b)\) holds for every
\(c\in J\). Take any \(z\) with \(a<z<d\).
If \(\psi_2\) is true at \(v=z\), the original matrix
holds at that point. If it is false, the alternative's universal
condition forces \(\rho_J^L(z,d)\).

When \(J=\varnothing\), that residual is false, so the latter
case is impossible. When \(J\ne\varnothing\), the finite
disjunction supplies an index \(c\in J\) with \(L_c(z,d)\).
The fixed-side conjunction supplies \(R_c(d,b)\), and the
reverse cut factorization yields \(\phi(z,b)\).
Thus every tested \(z\) satisfies the original matrix. The
same \(d\) witnesses the original existential formula.

No negative fixed-side assertion is required in this converse.
Some additional \(R_c(d,b)\) with \(c\notin J\) may be
true; that cannot invalidate an already supplied cut. This justifies
the source's passage from the complete truth patterns
\(\Omega_\alpha\) to their positive conjunctions
\(\Sigma_J^L\). It does not assert equivalence of each complete
truth pattern with its positive part in isolation.

\paragraph{Length, syntax and order scope.}
Every disjunct of \(\rho_J^L\) is one of the \(L_c\)
from the cut factorization. Its length is at most \(2n+1\),
and the only possible length-\(2n+1\) index is
\((\mathrm g,n)\). Because a subset lists each cut at most
once, there is at most one such disjunct. The empty residual has
no basic disjuncts at all, so it satisfies the bound as well.
All nonempty fixed-side conjuncts are basic patterns at \((u,y)\),
and all nonempty residual disjuncts are basic patterns at \((v,u)\).
Their elementary membership follows by the identity substitution
history. The temporary complete truth patterns are proof devices,
not an alteration of those positive output classes.

The outer strict witness remains in every alternative. If it does not
exist, both formulas are false. If the tested interval \((a,d)\)
is empty, the universal matrices impose no condition, and the empty
alternative supplies the required tautological fixed side. No density
is assumed. All equivalences concerned arbitrary structures and
assignments over \(\mathsf{Lin}\), so they have precisely that
scope. No property \(I\) over complete orders was invoked.

\paragraph{Verification of the two-interval example.}
In the body example the only points satisfying \(0<v<3\) are
1 and 2. The sequence \((1,2,6)\) has its point labels
\(A_0(1),A_1(2),A_2(6)\). Its first open gap is empty and
its second is \(\{3,4,5\}\subseteq B_2\). Hence it realizes
the original word at \((1,6)\). The sequence \((2,4,6)\)
has the required point labels; its first gap is
\(\{3\}\subseteq B_1\) and its second is
\(\{5\}\subseteq B_2\). Thus the original universal condition
holds with \(u=3\) and false \(\psi_2\).

On the fixed side \((3,6)\), the first-gap suffix is realized
by \((3,4,6)\). Its first point requires \(B_1(3)\), its
first gap is empty, its middle point requires \(A_1(4)\), its
last gap \(\{5\}\) satisfies \(B_2\), and its last point
requires \(A_2(6)\). The second-gap suffix is realized by
\((3,6)\), because \(B_2(3)\),
\(\{4,5\}\subseteq B_2\), and \(A_2(6)\) hold.
The point-cut suffix fails because its first point would require
\(A_1(3)\), which is false. The fixed-side truth subset is
therefore exactly \(\{(\mathrm g,1),(\mathrm g,2)\}\).

At \((1,3)\), the second-gap prefix is realized by \((1,2,3)\):
its point labels are \(A_0,A_1,B_2\) at those points, and
both intervening gaps are empty. The first-gap prefix fails there
because its open gap contains 2, whereas \(2\notin B_1\).
At \((2,3)\), the first-gap prefix holds: \(A_0(2)\) and
\(B_1(3)\) hold and the gap is empty. The second-gap prefix
would require a strict intermediate witness between 2 and 3, which
does not exist. The point-cut prefix fails at either pair because
its final point would require \(A_1(3)\).
Consequently neither singleton gap subset supplies both universal
tests, while their two-member subset does. The example verifies
the necessity of permitting different residual disjuncts at different
values of the quantified endpoint. Its cut list and lengths follow
by direct substitution of \(n=2\) in the word definitions.
\end{proof}
\end{annexitem}

\begin{annexitem}{annex:p9.kamp.iii2.lemma7.right}{thm:p09-kamp-iii2-lemma7-right}{The complete right-interior construction and its empty case}
The source's second orientation places the variable endpoint to the
right of the fixed cut. The proof below constructs its fixed-side
conjunctions and residual disjunctions explicitly, including the
empty aggregate needed when the other universal disjunct suffices
\cite[ch.~III, lemma~7(b), p.~81]{Kamp1968}.

\begin{proof}
Fix an arbitrary nonempty strict linear structure \(\mathcal M\),
the predicate interpretation, and an assignment to all free parameters.
Let \(a,b\) be the values of \(x,y\). Now the pattern is
\(\phi(x,v)\), and its word remains \(W\).
For each truth assignment \(\alpha:\mathcal C_n\to\{0,1\}\),
form the complete fixed-side condition
\[
 \widetilde\Omega_\alpha(x,u)\defeq
 \bigwedge_{c\in\mathcal C_n}
 \begin{cases}
 L_c(x,u),&\alpha(c)=1,\\
 \neg L_c(x,u),&\alpha(c)=0.
 \end{cases}
\]
For any values of \(x,u\), exactly one of these holds: its
indices assigned 1 are exactly those for which \(L_c\) is true.
The same finite truth-value argument applies at each index and yields
\(\bigvee_\alpha\widetilde\Omega_\alpha\) valid.
These conditions have no free \(v\), so the selected truth
pattern stays fixed while \(v\) ranges to the right of the cut.

\paragraph{Forward implication with the fixed prefix truths.}
Suppose the original right-interior formula holds. Choose its witness
\(d\) for \(u\). We have \(a<d<b\), the truth of
\(\psi_1\), and, for every \(z\) with \(d<z<b\),
the truth of \(\psi_2\) at \(v=z\) or the truth of
\(\phi(a,z)\). Put
\[
                    J=\{c\in\mathcal C_n:L_c(a,d)\text{ holds}\}.
\]
Every member gives its conjunct in \(\Sigma_J^R(a,d)\).
If there is no such member, \(\Sigma_\varnothing^R=T^K_u\)
holds by its definition.

Take an arbitrary \(z\) with \(d<z<b\). If \(\psi_2\)
is true there, the new universal matrix holds. Otherwise the original
matrix supplies \(\phi(a,z)\). Since \(a<d<z\), the
factorization lemma supplies \(c\) such that
\(L_c(a,d)\land R_c(d,z)\). Its first conjunct puts \(c\)
in the already fixed subset \(J\); its second gives the
corresponding disjunct of \(\rho_J^R(d,z)\). Thus the new
matrix holds at every tested point. With the unchanged outer
inequalities, unchanged \(\psi_1\), and the same witness
\(d\), this proves the alternative indexed by \(J\).
There is no choice of one cut required to work for every \(z\).

\paragraph{Reverse implication from each aggregate.}
Suppose some right output alternative is true, with subset \(J\)
and witness \(d\). Its outer conditions supply \(a<d<b\)
and \(\psi_1\). Its conjunction supplies \(L_c(a,d)\)
for every \(c\in J\). For any \(z\) with \(d<z<b\),
if \(\psi_2\) is true there, the original matrix follows.
If \(\psi_2\) is false, the new matrix gives
\(\rho_J^R(d,z)\).

For \(J=\varnothing\) this is impossible because the residual
is \(F^K_v\). For nonempty \(J\), one of its disjuncts
supplies \(R_c(d,z)\) for some \(c\in J\). Combining
it with \(L_c(a,d)\) and using \(a<d<z\), the reverse
factorization gives \(\phi(a,z)\). Hence the original
universal matrix holds at every point in its guarded interval,
and the same witness \(d\) proves the original formula.
As in the left case, the converse uses only the positive fixed-side
truths, so additional true prefixes outside \(J\) are harmless.

\paragraph{The unchanged original bound in the right orientation.}
Every residual disjunct is an \(R_c\). A point cut has right
length \(2(n-k)+1\), strictly below \(2n+1\), and a gap cut
has length \(2(n-k+1)+1\), at most \(2n+1\).
The sole equality case is the first gap, \((\mathrm g,1)\).
Each subset includes that index at most once. Thus every right
residual contains at most one disjunct of the original length and
no longer one. The empty residual contains no basic disjuncts.
Each nonempty fixed-side conjunct lies in
\(\operatorname{Bp}(x,u)\), and each nonempty residual
disjunct in \(\operatorname{Bp}(u,v)\), exactly as their
word constructions show. All operations are finite.

\paragraph{Quantifier scopes and all empty cases.}
The input formulas \(\psi_1,\psi_2\) were arbitrary. Neither
was replaced, simplified or assumed independent of \(u\).
The outer \(\exists u\) binds their free \(u\) occurrences
within its scope, as it did initially. The inner \(\forall v\)
binds the occurrences in its own matrix and leaves any free
\(v\) of \(\psi_1\) outside that matrix unchanged.
Fresh internal bound names in every cut pattern prevent additional
capture. All further parameters retain the initial assignment.
Consequently the two directions prove formula equivalence, not merely
an equivalence after restricting the other formulas to unary predicates.

If no \(d\) satisfies \(a<d<b\), both the original and every
output alternative are false. If such a \(d\) exists but
\((d,b)\) is empty, the original universal matrix is true
independently of its contents, and the empty alternative imposes the
same vacuous matrix with a tautological fixed side. If the interval
is nonempty and \(\psi_2\) holds throughout it, the same empty
alternative works with its false residual. The construction thus
retains the cases where no fixed-side basic pattern is available.

\paragraph{Verification of the all-false and smallest-word example.}
In the three-point order \(0<1<2\), with every unary predicate
false, \(P_0(w)\lor\neg P_0(w)\) is true at every point and
\(P_0(w)\land\neg P_0(w)\) is false there. With \(x=0,y=2\),
choose \(u=1\). The outer order condition and \(\psi_1=T^K_u\)
hold in both orientations. The corresponding universal condition
holds because \(\psi_2=T^K_v\) is true at every tested point;
in this three-point instance both tested open intervals are also empty.
Their empty-subset alternatives have exactly the same witness,
true fixed-side aggregate, and true universal implication.

Every basic past pattern begins with a unary point atom at its initial
endpoint, both at stage one and in each higher-stage constructor.
That atom is false under this interpretation, so every basic pattern
at every ordered endpoint pair is false. A nonempty conjunction of
such patterns is false. This verifies why imposing one basic pattern
as the fixed-side condition would remove a valid case. The same
argument works on \(0<1<2<3\) with \(x=0,y=3,u=1\) in the
right orientation, where the tested interval contains 2 and its
truth comes from \(\psi_2\), rather than from vacuity.

At \(n=1\) there is one gap cut, so the complete subset family
has two members. Substituting the word \((A,B,C)\) into the
cut definitions gives \((A,B,B)\) and \((B,B,C)\), each
with one interval and length three. The cut factorization proves
their conjunction equivalent to the original pattern under a strict
interior cut. Both directions above apply with these exact words;
neither uses a fictitious zero-interval pattern.

Every step used only finite word operations, ordinary Boolean logic,
the original quantifier scopes, and the linear-order cut equivalence.
Hence the right reduction, as well as the left, holds throughout
\(\mathsf{Lin}\), and its restriction to \(\mathsf{Com}\)
is available for the subsequent interval-property induction.
\end{proof}
\end{annexitem}
